\documentclass[11pt,a4paper]{article}
\usepackage[T1]{fontenc}
\usepackage[landscape,margin=12mm]{geometry}
\usepackage{graphicx,pgffor,hyperref}
\setlength{\parindent}{0pt}
\newcommand{\AtlasDataset}[3]{%
 \foreach \level [count=\groupindex from 1] in {#2}{%
  \edef\groupid{\ifnum\groupindex<10 0\fi\groupindex}%
  \foreach \method in {#3}{%
   \clearpage
   \section*{Dataset #1, \method, group \groupindex}
   \includegraphics[width=\linewidth,height=.80\textheight,keepaspectratio]{figures/atlas/set#1_group\groupid_\method.png}
   \par\small Blue: reference; orange: reconstruction. Signed errors in mVs.
   White regions have no common numerical support. Dataset 4 is a processed temperature fixture.
  }%
 }%
}
\title{Coefficient-Based Flux Maps: Complete Group Atlas}
\author{Flux Map Research Study}\date{24 October 2026}
\begin{document}\maketitle
Both stator flux components are shown for every group and applicable model.
Figure titles identify training, validation-refit, independent excitation test and upper-boundary cases.
Upper-boundary reconstruction is unavailable outside the physical training range.
\par\medskip
The d-flux 3D view exchanges displayed d/q axes. Error maps retain physical d/q order.
Scales are shared within a dataset. Reference surfaces preserve source-triangle gaps;
signed errors use common support across methods. Display grids do not replace point-based metrics.
The intersection can also mask finite coefficient-model predictions, including training
boundary slices where the raster LUT has no complete interpolation cells.
\par\medskip
Dataset 4 has unverified acquisition provenance and two temperature bands. It is a processed-reference
implementation exercise with spatial holdout, not independent temperature-validation evidence.
\centering\AtlasDataset{1}{2,5,8,11,14,17,20,23,26,29,32,36,40}{RBF,Spline,LUT,kA,kB,kC,kD,kE,kF,kG}
\AtlasDataset{2}{0,2,5,8,11,14,17,20,23,26,29,32,36,40}{RBF,Spline,LUT,kA,kB,kC,kD,kE,kF,kG}
\AtlasDataset{3}{3,6,9,12,15,18,21,24,27,30}{RBF,Spline,LUT,kA,kB,kC,kD,kE,kF,kG}
\AtlasDataset{4}{30,70}{RBF,Spline,LUT}

\end{document}
